\lecture{12}
\title{Discrete Fourier Transform I}

\begin{document}

\subsection{Lecture: Discrete Fourier Transform}

The problem with Fourier transforms is that, unlike Fourier series,
they are not easily expressible in code:
they're continuous data, not discrete.
The \textbf{discrete Fourier transform}, as its name suggests,
is the fix for this.

\textbf{Definition.}
Given a possibly aperiodic DT signal $x[n]$
and any positive integer $N$,
the \textbf{discrete Fourier transform (DFT)} of $x[n]$
is the same as the Fourier series of the periodic signal
$x'[n] := x[n \% N]$.

\begin{remark}
    Recall that this means the DFT of $x[n]$ is an array $X$,
    where $X[k]$ denotes the coefficient of $e^{j(2\pi k/N)n}$
    in the DT Fourier series of $x'[n]$.
    Note that $X$ is an infinite periodic array with period $N$.
\end{remark}

In some ways, the DFT is the culmination of all of the Fourier analysis
in 6.300 up to this point: it combines
the signal-aperiodicity of the DTFT
with the computational feasibility of the DTFS.

\begin{center}
    \frame{\includegraphics[width=0.4\textwidth]{lecture-12/chart}}
\end{center}

\textbf{Example.}
The DFT of $x[n] = \cos\left(\frac{2\pi n}{64}\right)$ for $N = 64$
is the same as the Fourier series of $x[n]$, plotted below.
\begin{center}
    \frame{\includegraphics[width=0.4\textwidth]{lecture-12/periodic-DFT}}
\end{center}

\textbf{Example.}
The DFT of $x[n] = \cos\left(\frac{3\pi n}{64}\right)$ for $N = 64$,
plotted below, differs from the Fourier series of $x[n]$.
\begin{center}
    \frame{\includegraphics[width=0.4\textwidth]{lecture-12/aperiodic-DFT}}
\end{center}
The reason why is that $x[n]$ does not have period $N = 64$; its period is $N = 128$.
However, the DFT of $x[n]$ for $N = 64$ does match the Fourier series
of the signal pictured below.
\begin{center}
    \frame{\includegraphics[width=0.4\textwidth]{lecture-12/forced-periodic}}
\end{center}

Unsurprisingly, here's what the analysis and synthesis equations
for the DFT look like.

\begin{theorem}{DFT Formulas}
    Given a DT signal $x[n]$ and a positive integer $N$, the DFT of $x[n]$ is
    $X[k] := \frac{1}{N} \sum_{n = 0}^{N - 1} x[n] e^{-j \frac{2\pi k}{N} n}$,
    which satisfies
    $x[n] = \sum_{k = 0}^{N - 1} X[k] e^{j \frac{2\pi k}{N} n}$.
\end{theorem}

These formulas look \emph{exactly} the same as the DT Fourier series
formulas presented in \href{5.html#lecture-dt-fourier-series}{Lecture 5}, expectedly.

\subsection{Lecture: DFT vs. DTFT and Windows}

We know the relationship between the DFT and the DTFS:
they're exactly the same.  What about the DFT and the DTFT?

\begin{theorem}{DFT vs. DTFT}
    Consider a DT signal $x[n]$ and a positive integer $N$.
    Define a new signal $x_w[n]$ as follows:
    \[
        x_w[n] := \begin{cases}
        x[n] & \text{ if } n \in [0, N - 1]. \\
        0 & \text{ otherwise.}
        \end{cases}
    \]
    Suppose $X[k]$ is the DFT of $x[n]$,
    and $X_w(\Omega)$ is the DTFT of $x_w[n]$.  Then
    $X[k] = \frac{1}{N} X_w\left(\frac{2\pi k}{N}\right)$.

    In other words,
    $X[k]$ is the result of taking $N$ samples from $X_w$
    at an interval of $\frac{2\pi}{N}$
    and dividing all samples by $N$.
\end{theorem}

\begin{proof}
    Recall from
    \href{7.html#lecture-the-dt-fourier-transform}{Lecture 7}
    that the DTFT $X_w(\Omega)$ is defined via
    $X_w(\Omega) = \sum_{n = -\infty}^{\infty} x_w[n] e^{-j \Omega n}$.
    Comparing with the formula
    $X[k] := \frac{1}{N} \sum_{n = 0}^{N - 1} x[n] e^{-j \frac{2\pi k}{N} n}$
    for the DFT, the relationship is immediate. ~ $\blacksquare$
\end{proof}

The diagram below paraphrases the statement of the theorem above.

\begin{center}
    \frame{\includegraphics[width=0.4\textwidth]{lecture-12/DFT-versus-DTFT}}
\end{center}

Note that the constant $N$ decides the size of the ``window''
that our DFT takes into account.

\begin{problem}
    Compute the DFT of the pure-frequency DT signal $x[n] := e^{j\Omega_0 n}$
    for varying window sizes $N$.
\end{problem}

\begin{solution}
    Equivalently, we compute the DTFT of
    $w[n] \cdot x[n]$, where $w[n]$ is the signal given by:
    \[ w[n] := \begin{cases}
        1 & \text{ if } n \in [0, N - 1] \\
        0 & \text{ otherwise.}
    \end{cases} \]
    Note that the DTFT of $w[n]$ itself looks like the following:
    \[ W(\Omega) = \sum_{n = -\infty}^{\infty} w[n] e^{-j \Omega n}
    = \sum_{n = 0}^{N - 1} e^{-j \Omega n}
    = \frac{1 - e^{-j \Omega N}}{1 - e^{-j \Omega}}
    = \left[\frac{
        \sin \left(\frac{N \Omega}{2}\right)
    }{
        \sin \left(\frac{\Omega}{2}\right)
    }\right] \cdot e^{-j \Omega(N - 1) / 2}. \]
    Now we invoke a ``dual'' of one of the properties from
    \href{
        7.html#lecture-and-recitation-properties-of-the-dt-fourier-transform
    }{Lecture 7}: the DTFT of $w[n] \cdot x[n] = w[n] \cdot e^{j\Omega_0 n}$
    is just $W(\Omega - \Omega_0)$.

    Thus, the DFT of $x[n]$ is the result of
    sampling from $W(\Omega - \Omega_0)$
    at intervals of $\frac{2\pi}{N}$
    and dividing the samples by $N$.
    \begin{center}
        \frame{\includegraphics[width=0.4\textwidth]{lecture-12/clean-DFT}}
    \end{center}
    For $\Omega_0 = \frac{4\pi}{15}$ (shown above),
    the resulting DFT reads as pretty strong evidence
    that the original DT frequency was $\frac{4\pi}{15}$.
    \begin{center}
        \frame{\includegraphics[width=0.4\textwidth]{lecture-12/ugly-DFT}}
    \end{center}
    For $\Omega_0 = \frac{3.2\pi}{15}$ (shown above),
    the resulting DFT is less informative\dots

    We can get more informative DFTs
    by increasing the window size $N$.
    The larger the value of $N$,
    the finer our frequency resolution,
    and the more samples we get to read.
    \begin{center}
        \frame{\includegraphics[width=0.3\textwidth]{lecture-12/finer-windows}}
    \end{center}
\end{solution}

\subsection{Recitation: Spectral Analysis (Interpreting the DFT)}

\begin{problem}
    A DT signal $x[n]$ is produced by sampling from a CT signal
    with sampling frequency $f_s$,
    and its DFT $X[k]$ is computed using a window length of $N$.
    Assume $x[n]$ is periodic modulo $N$.
    What does $X[7]$ say about the original CT signal?
\end{problem}

\begin{solution}
    Since $x[n]$ is periodic modulo $N$,
    the DFT of $x[n]$ is very informative:
    it's exactly the same as the DTFS of $x[n]$.

    The rest is dimensional analysis.
    Note that $X[7]$ is the coefficient of $e^{j(2\pi(7)/N)n}$
    in the DTFS, so it represents the strength of the
    discrete-time angular frequency $\Omega_0 = 7 \times \frac{2\pi}{N}$.
    
    In continuous-time, this translates to
    (recalling \href{4.html#recitation-periods-of-dt-signals}{Lecture 4})
    a continuous-time angular frequency of
    $\omega_0 = 7 \times \frac{2\pi}{N} \times f_s$.
\end{solution}

More generally, we can say the following:

\begin{theorem}{Spectral Analysis}
    The DFT $X[k]$ measures the strength of the
    continuous-time cyclical frequency $f_0 = \frac{k}{N} \times f_s$.
\end{theorem}

\begin{problem}
    Suppose $f_s = 44 \, 100$ and a DFT $X[k]$ using window size $N = 100$.
    Given that $X[13] = X[87] = 1$ and all other entries of $X$ are zero,
    what can we say about the original CT signal?
\end{problem}

\begin{solution}
    The original CT signal probably had a cyclical frequency
    of $\frac{13}{100} \times 44 \, 100 = 5 \, 733$ Hz.
\end{solution}

Notably, the data provided by the DFT is discrete,
so the set of CT frequencies it may tell us about is also discrete.

\textbf{Definition.} Two angular frequencies $\Omega_1$ and $\Omega_2$
are \textbf{resolved} (that is, distinguishable in the DFT)
if $|\Omega_1 - \Omega_2| > \Delta\Omega$,
where $\Delta\Omega := \frac{2\pi}{N}$ is the
\textbf{frequency resolution.}
(For cyclical frequencies, it looks like $\Delta f := \frac{f_s}{N}$.)

\end{document}